\documentclass[12pt]{article}
\usepackage[T1]{fontenc}
\usepackage[utf8]{inputenc}
\usepackage{lmodern}
\usepackage{textcomp}
\usepackage{enumitem}
\usepackage{geometry}
\usepackage{graphicx}
\usepackage{amsmath, amssymb}
\geometry{margin=1in}
\begin{document}
\begin{center}
History - K-12 Level Assessment 24 \
Total Marks: 40 \
Answer all questions.
\end{center}

\section*{Multiple Choice (10 marks)}
\begin{enumerate}[label=\arabic*.]
\item Unlike Mesopotamian writings, the records left by the Maya are principally concerned with:
\begin{enumerate}[label=(\alph*)]
    \item political and military history.
    \item philosophy and religion.
    \item trade and economics.
    \item astronomy and mathematics.
\end{enumerate}
\item The idea that cultures change through time in relation to their environments, along many different paths, is known as:
\begin{enumerate}[label=(\alph*)]
    \item unilineal evolution.
    \item environmental uniformitarianism.
    \item multicultural adaptation.
    \item multilineal evolution.
\end{enumerate}
\item Which foods are included in the triad of plants that provided the subsistence base for indigenous New World civilizations?
\begin{enumerate}[label=(\alph*)]
    \item maize, beans, and squash
    \item maize, lentils, and wheat
    \item rice, beans, and squash
    \item yams, maize, and millet
\end{enumerate}
\item What was the unifying element of Olmec culture?
\begin{enumerate}[label=(\alph*)]
    \item religious iconography
    \item large armies
    \item agricultural surpluses
    \item all of the above
\end{enumerate}
\item The ridge of bone that runs along the top of the skull in Australopithecus robustus is called a:
\begin{enumerate}[label=(\alph*)]
    \item diastema.
    \item cranium.
    \item maxilla.
    \item sagittal crest.
\end{enumerate}
\end{enumerate}

\section*{True / False (10 marks)}
\textit{Answer True or False}
\begin{enumerate}[label=\arabic*.]
\setcounter{enumi}{5}
\item True or False: Excavation of Yang-shao sites in China reveals the presence of wheat, rice, and peas as domesticated crops.
\item True or False: Primatology refers to the study of the Earth's layers and geological formations.
\item True or False: Pompeii's preservation was due to avalanches pushing the remains into the sea, where seawater helped to preserve them.
\item True or False: Capacocha was an Incan practice of feasting on the hearts of slain enemies to honor their bravery in battle.
\item True or False: The pharaoh ruled a population of 1,000,000 people, of which 50\% were farmers.
\end{enumerate}

\section*{Long Form Response (20 marks)}
\textit{Answer each question with a clear and complete explanation.}
\begin{enumerate}[label=\arabic*.]
\setcounter{enumi}{10}
\item Discuss the key differences between Aurignacian and Mousterian technologies, focusing on the advancements in tool design and their implications for early human life.
\item Discuss the significance of the appearance of early human ancestors in the context of the universe's history as represented in the author's film. Why is this timing important?
\end{enumerate}

\vfill
\noindent\textit{End of Paper}
\end{document}